The measured PDD and beam-profile data for photon, FFF photon, and electron beams were consistent with the expected clinical behavior of a medical linear accelerator.

For flattened photon beams, the profile at smaller field size remained comparatively uniform in the central region, whereas the larger field size showed a clearer horn effect near the off-axis region. This is physically expected because the flattening filter is designed to compensate central-axis fluence by hardening and attenuating the beam more at the center than at the periphery. At larger field widths, this over-compensation becomes more apparent and produces peripheral dose elevation (horns). In clinical dosimetry, this confirms why profile evaluation at large fields is essential for baseline QA, monitor-unit confidence, and TPS beam-model validation.

The FFF beam profile showed a non-flat, forward-peaked distribution with a higher central axis dose and gradual off-axis fall-off. This is the characteristic signature of filter-free operation: after removing the flattening filter, there is reduced head scatter and less beam hardening from the filter, so lateral fluence compensation is absent. From a medical physics perspective, this behavior is expected and acceptable, provided machine-specific baseline quantities (inflection points, unflatness metrics, and penumbra definitions for FFF mode) are stable and reproducible for periodic QA.

The electron PDD exhibited a high dose in the near-surface to shallow-depth region, followed by a relatively rapid distal fall-off and a low-dose bremsstrahlung tail. This reflects the finite range of therapeutic electrons in water-equivalent media: electrons deposit dose efficiently up to a practical range and then lose dose sharply due to energy depletion and multiple scattering. Clinically, this supports the use of electrons for superficial targets while reducing deeper normal-tissue dose compared with photons.

Overall, the observed depth-dose and profile characteristics were in line with standard radiotherapy beam physics. The extracted dosimetric parameters (field width behavior, symmetry/flatness-type checks, penumbra trends, and depth-dose pattern) indicate that the beams are suitable for clinical implementation, subject to routine constancy checks and continued TPS-measurement agreement.
